\subsection{Three visible cuts in the seven-deletion case}
This subsection excludes the branch in which all three deleted additions enter the quadratic frame directly, so all three cuts are visible in the sense defined at the start of Section~\ref{sec:four}.

\begin{proposition}[Three visible cuts]\label{prop:threevisible}
Suppose exactly four quadratic witnesses are missing, $d=4$, the extracted quadratic circuit is optimal, and seven deleted gates consist of four products and three additions. The three additions cannot all be visible cuts.
\end{proposition}
\begin{proof}
Apply full-carrier erasure before restricting variables. Every private gate must remain: an erasure or deletion of an unused private gate would give at most six deletions without changing the extracted frame, contrary to the six-deletion exclusion. Choose the assigned minimum cover $Z$ of $H$. After setting its variables to zero, the graph on $T=V(G)\setminus Z$ is the four-edge matching with isolated vertices. Original private gates that become zero or redundant on this restriction stay in the count.

Each cut's high operand is a private-product output. To see this, first note that an old degree-one value is too small to be a high operand. A visible cut of minimum degree one is likewise too small. A selected cut of minimum degree two has an addition-only path to the quadratic output, since any product on that path would raise the minimum degree. Its full value is therefore a protected quadratic carrier. Full-carrier erasure removes every off-carrier use of such a value; it cannot supply the high operand of a degree-one cut. It cannot supply the high operand of a degree-two cut either, because that operand has degree strictly greater than two. The same reasoning excludes old quadratic values. All three deleted additions are already the visible cuts, so no other kind of high operand remains.

Write $p_1,\ldots,p_4$ for the private products in topological order, and mark those selected as high operands. Every operand of a private product is either an earlier private-product output or an old positive linear clause OR a subset of earlier marked outputs. Indeed, between private products only the old degree-one additions and the three cut insertions remain. Induction through those additions gives this form. An earlier unmarked product cannot be added to an old clause without another deleted addition. Carrier inputs have already been erased. Making each marked output available immediately after its creation enlarges the possible circuits and includes the original one.

Temporarily replace each high-input occurrence by an independent symbol of weight equal to the minimum degree of its low operand. The exact quadratic frame has weighted degree two. A weight-two symbol has only a scalar context; a weight-one symbol has a linear context or a second weight-one symbol. The square of one symbol cannot occur, since substituting its nonzero old linear low part would give a square in the squarefree quadratic polynomial, and positive coefficients cannot cancel it. Thus, on the restriction, the output has the form
\begin{equation}\label{eq:threevisibleform}
 \bigvee_{i=1}^3 g_iL_i\ \vee\!
 \bigvee_{1\leq i<j\leq3:\,\delta_{ij}=1}g_ig_j,
 \qquad \delta_{ij}\in\{0,1\},
\end{equation}
where $g_i$ are the marked products, repetitions allowed, and each $L_i$ is an affine positive clause. Allowing all pair terms is harmless even when particular weights would forbid some of them.

Every nonzero nonconstant $L_i$ contains an isolated variable or both endpoints of a matching edge. In fact, such a direct context comes from perturbing an assigned center $h\in Z$ and contains the nonempty set $S=N_H(h)\cap T$. If the context contained no isolate and no complete matching pair, then $S$ would meet exactly $|S|$ different matching edges. Replace $h$ in $Z$ by $S$, and add one endpoint of each remaining missing edge. This covers $H$: old edges from $h$ into $T$ meet $S$, and old edges between $h$ and another center meet the other center. It also covers the four missing edges, with at most $|Z|-1+|S|+4-|S|=|Z|+3$ vertices. This contradicts $d=4$.

If the three cuts select at most two distinct product nodes, collect the terms in \eqref{eq:threevisibleform} with the same high value. A repeated square is Booleanly that value and hence has constant context one. Unions of permitted contexts are still permitted. The resulting function belongs to the two-high-signal class excluded in Proposition~\ref{prop:fiveproductseven}, whose proof permits fewer than five products and does not require the restricted signals to be nonzero or distinct. Thus it remains to consider three distinct marked products.

Every original private product is ancestral to a marked product. Its first entry into the selected subgraph must be along the high-input edge of a visible cut: entry through a product, an equal-degree sum, or a cut's low input would make the private node selected already. The last of the four topologically ordered products must therefore itself be marked. Consequently the three marked index sets are exactly
\[
 \{1,2,4\},\qquad \{1,3,4\},\qquad \{2,3,4\}.
\]
Keep these original indices even if two restricted values coincide or vanish.

For completeness, the finite obstruction for these three classes is specified as follows. Identify all isolated variables with one variable $z$, leaving the target matching unchanged. In each product operand select exactly one of two kinds of input: an arbitrary old linear clause in the eight endpoints and $z$, optionally joined by any earlier marked outputs; or one earlier private-product output alone. The empty clause is allowed. Select each affine context by its nine variable-membership bits and a constant-one bit. Require a nonzero nonconstant context to contain $z$ or a full matching pair, and select the three pair terms independently. No essentiality, nonvanishing, or distinction of restricted truth functions is imposed.

The circuit must be true on each of the four edge assignments with $z=0$ and false on each of the sixteen maximal independent endpoint transversals with $z=1$. These twenty conditions imply equality on every assignment by monotonicity: every edge-containing assignment dominates an edge test, and every independent assignment is below a tested transversal. Exact Tseitin equivalences express the specified evaluations. The resulting CNFs have, respectively,
\[
 (3342,11778),\qquad(3300,11618),\qquad(3258,11458)
\]
variables and clauses. Complete DRUP proofs for them contain $196\,938$, $236\,242$, and $286\,925$ proof lines. Independent regeneration reproduced every CNF byte for byte. DRAT-trim verified each refutation in RUP-only mode with zero RAT steps, in $6.20$, $7.47$, and $10.62$ seconds. The encoder \emph{Visible-cut obstruction}, compressed refutations, hashes, regeneration checks and checker logs are available through the workbench. Reverse unit propagation establishes each added clause from earlier clauses and ends in the empty clause, as in Lemma~\ref{lem:threeprivatefinite}. These three refutations exclude every remaining class and prove the proposition.
\end{proof}

The remaining visibility cases have one or two visible cuts, considered next.
